% Section 9 -- Conclusion. Plain, direct prose (RAG register): the conceptual contribution, the
% mechanism in relation to DCPO / Bereket & Leskovec, the triangulation as a diagnostic, limitations,
% future work. No editorial teasers, no clefts.
\section{Conclusion}
\label{sec:conclusion}

We have shown that reinforcement learning can calibrate a probabilistic forecaster to the level of the
betting market using only realized game outcomes, without human labels or supervised fine-tuning.
Verifiable-reward training has used the realized label as the reward, which is well behaved when that
label is deterministic but injects irreducible variance when the outcome is stochastic. We reward an
aggregate of those same labels, a state-conditioned empirical rate, which is equally verifiable and far
less noisy, and which turns a proper scoring rule into a usable training signal for events that have no
correct answer. This carries verifiable-reward training from tasks with a right answer to forecasting,
where the answer is itself a probability.

Even with the denoised target, taking the gradient through a chain of thought decalibrates the model,
because optimizing the final probability rewrites the reasoning that produced it. This refines earlier
accounts: the accuracy-calibration gradient conflict \citep{dcpo} and the overconfidence attributed to
the group-standard-deviation normalization of the advantage \citep{uncalibrated_reasoning} both take the
realized outcome as the target, whereas with a denoised target that normalization is instead necessary
and the problem that remains is the gradient's reach into the reasoning. Confining the gradient to the answer
removes it, by dropping the reasoning or by masking the gradient on it, and the masked variant keeps a
chain of thought from which the forecast follows.

Three unrelated estimators converge on the same Brier score: a reinforcement-trained small model, a
frontier model, and a tabular rate. The convergence is strong evidence that the score is the limit set
by the information in the public game state, not a property of any one of them. It is also a practical
test of whether a forecaster has exhausted its inputs: once unrelated methods agree, further accuracy
must come from new information, not a larger model. A market, where one exists, supplies the
independent reference that makes the test sharp.

The method reaches a strong reference only where the public state already carries most of the predictive
signal, and where outcomes are dense and resolved quickly enough to estimate a reliable empirical rate;
sparse or long-horizon events would weaken both the reward and the comparison against a market. Within
these limits the empirical rate could be enriched with more state features or replaced by a learned rate
model to raise the ceiling it imposes, and the same reward and gradient mask should extend to other
aleatoric domains where calibrated probabilities matter and outcomes are eventually observed, from
weather to elections to clinical prognosis. Calibrated forecasting can be learned from outcomes alone,
and a small open model can match a market once its reward and its gradient are aligned with the quantity
it is asked to forecast, while still exposing the reasoning behind each prediction.
